% El text de la pauta s'ha transcrit amb les errates d'accentuació de l'original («Es un so», «esta format», «es el pic A», «te una»).
\begin{solapartat}
\punts{0,3} Es un so complex ja que l'espectre mostra que esta format per la superposició de sons de moltes freqüències diferents.

\medskip
\punts{0,4} El pic que correspon a la freqüència fonamental es el \textbf{pic A}, que deu tenir una freqüència, comparant-lo amb el pic B que es el segon harmònic, de $\dfrac{880}{2} = \mathbf{440}$ \textbf{Hz}.

\medskip
\punts{0,3} Els diferents pics tenen freqüències f\textsubscript{n} = n·f\textsubscript{0} =n·440 s\textsuperscript{-1}. El pic I seria el 9è, així que $f_9 = 9\cdot 440 =$ \textbf{3960 Hz.}
\end{solapartat}

\begin{solapartat}
% A l'original la gràfica és a la dreta dels càlculs d'aquest apartat.
\figura[0.4]{sol_fig1.png}

\punts{0,2} $L = 10\log\dfrac{I}{I_0}$

\smallskip
\punts{0,2} $I = I_0 10^{\frac{L}{10}}$

\smallskip
\punts{0,4}

\smallskip
$\left.\begin{aligned} I_C &= I_0\cdot 10^{8,7}\\ I_F &= I_0\cdot 10^{6} \end{aligned}\right\} \rightarrow \dfrac{I_C}{I_F} = \dfrac{I_0\cdot 10^{8,7}}{I_0\cdot 10^{6}} = 501,2 \approx 500$

\medskip
\punts{0,2} El pic C te una intensitat sonora aproximadament 500 vegades major que el pic F.
\end{solapartat}
